% !TeX program = xelatex
\documentclass[UTF8,a4paper,12pt]{ctexart}
\usepackage{geometry}
\usepackage{listings}
\usepackage{xcolor}

\geometry{a4paper, top=0.4in, bottom=1in, left=1in, right=1in}

\lstset{
    language=C++,
    basicstyle=\ttfamily\small,
    keywordstyle=\color{blue},
    commentstyle=\color{gray},
    stringstyle=\color{orange},
    numbers=left,
    numberstyle=\tiny\color{gray},
    frame=single,
    breaklines=true
}

\title{取余运算的应用}
\author{GESP 2级}
\date{}

\begin{document}

\maketitle

\section*{取余运算的应用 \\ (Modulo Operation)}

\begin{lstlisting}
// 判断奇偶
if (n % 2 == 0) cout << "偶数";
else cout << "奇数";

// 取个位、十位、百位
int num = 123;
int ones = num % 10;            // 3（个位）
int tens = num / 10 % 10;      // 2（十位）
int hundreds = num / 100 % 10;     // 1（百位）

// 判断整除
if (a % b == 0) cout << "a能被b整除";

// 时钟问题：现在是10点，过5小时后是几点？
int current = 10;
int hours = 5;
int result = (current + hours) % 12;  // 3点（下午）
int hour24 = (current + hours) % 24;   //24小时制 15点
\end{lstlisting}

\section*{词汇表 (Glossary)}

\begin{center}
\begin{tabular}{|l|l|p{10cm}|}
\hline
\textbf{英文名词} & \textbf{中文名词} & \textbf{解释} \\
\hline
Modulo & 取余 (\%) & 求两个整数相除后的余数。如 10 \% 3 = 1，余数为1。 \\
\hline
Parity Check & 奇偶判断 & 通过 \texttt{n \% 2} 判断奇偶，余数为0是偶数，为1是奇数。 \\
\hline
Digit Extraction & 数位提取 & 通过 \texttt{num \% 10} 取个位，\texttt{num / 10 \% 10} 取十位，\texttt{num / 100 \% 10} 取百位。 \\
\hline
Divisibility & 整除 & 当 \texttt{a \% b == 0} 时，称 a 能被 b 整除，余数为0。 \\
\hline
\end{tabular}
\end{center}

\end{document}
